\documentclass[11pt]{article}
\usepackage[utf8]{inputenc}
\usepackage[T1]{fontenc}
\usepackage{booktabs}
\usepackage{amsmath}
\usepackage{geometry}
\geometry{margin=2.5cm}

\title{Causal Fairness Analysis Report}
\author{Generated by FairMind and an LLM}
\date{2026-08-19}

\begin{document}
\maketitle

\section{Analysis setup}

\begin{tabular}{ll}
\toprule
Dataset & adult \\
Protected attribute ($X$) & sex \\
Comparison & Female $\to$ Male \\
Outcome ($Y$) & income \\
Mediator ($W$) & hours-per-week \\
Confounder ($Z$) & education \\
\bottomrule
\end{tabular}

\section{Causal effects (FairMind, exact values)}

The five effects below are computed by FairMind through exact inference on the
fitted Bayesian network. These values are final and must not be recomputed or
altered.

\begin{tabular}{lr}
\toprule
Effect & Value \\
\midrule
Total Variation (TV) & 0.193714 \\
Total Effect (TE) & 0.184736 \\
Spurious Effect (SE) & 0.008977 \\
Direct Effect (DE) & 0.138404 \\
Indirect Effect (IE, reverse form) & -0.046333 \\
\bottomrule
\end{tabular}

\section{Qualitative interpretation}

\subsection{Total effect and spurious component (TV, TE, SE)}

The total disparity between females and males in income is 0.193714. After accounting for the confounding variable education, the genuine causal effect (TE) is 0.184736, indicating a significant portion of the observed disparity is real. The spurious effect (SE) of 0.008977 is relatively small, suggesting that the majority of the observed difference is driven by the direct effect of sex rather than by the confounder education.

\subsection{Direct effect (DE)}

The direct effect (DE) of 0.138404 indicates a substantial direct discrimination against females, as the income disparity persists even when the mediator (hours-per-week) is held constant. This direct effect is notably larger than the total effect (TE), suggesting that the direct discrimination is a major contributor to the overall disparity.

\subsection{Indirect effect (IE)}

The indirect effect (IE) of -0.046333 shows that the mediator (hours-per-week) offsets the direct effect, reducing the overall disparity. This suggests that the channel through the mediator mitigates the direct discrimination against females, as the indirect effect and direct effect have opposite signs.

\section{Recap Questions}

\begin{enumerate}
\item Is there direct discrimination against the protected group? \textbf{LLM:} YES \quad \textbf{Expected:} YES \quad (correct)
\item Does the indirect effect through the mediator mitigate the total effect? \textbf{LLM:} YES \quad \textbf{Expected:} YES \quad (correct)
\item Does the observed total effect exceed the practical relevance threshold? \textbf{LLM:} YES \quad \textbf{Expected:} YES \quad (correct)
\item Does the spurious component (SE) contribute substantially to the observed difference? \textbf{LLM:} NO \quad \textbf{Expected:} NO \quad (correct)
\item Is the direct effect the dominant component compared to the indirect effect? \textbf{LLM:} YES \quad \textbf{Expected:} YES \quad (correct)
\end{enumerate}


\vspace{1em}
\noindent\textbf{Answer consistency:} 5/5 (100.0\%). The expected answers are derived deterministically from the FairMind values alone, through threshold rules, and were added to the document after it was generated: the model never saw them.

\end{document}